%%%%%%%%%%%%%%%%%%%%%%%%%%%%%%%%%%%%%%%%%
% Journal Article
% LaTeX Template
% Version 2.0 (February 7, 2023)
%
% This template originates from:
% https://www.LaTeXTemplates.com
%
% Author:
% Vel (vel@latextemplates.com)
%
% License:
% CC BY-NC-SA 4.0 (https://creativecommons.org/licenses/by-nc-sa/4.0/)
%
% NOTE: The bibliography needs to be compiled using the biber engine.
%
%%%%%%%%%%%%%%%%%%%%%%%%%%%%%%%%%%%%%%%%%

%----------------------------------------------------------------------------------------
%	PACKAGES AND OTHER DOCUMENT CONFIGURATIONS
%----------------------------------------------------------------------------------------

\documentclass[
	a4paper, % Paper size, use either a4paper or letterpaper
	10pt, % Default font size, can also use 11pt or 12pt, although this is not recommended
	unnumberedsections, % Comment to enable section numbering
	twoside, % Two side traditional mode where headers and footers change between odd and even pages, comment this option to make them fixed
]{LTJournalArticle}

%\addbibresource{paramoji.bib} % BibLaTeX bibliography file
\bibliographystyle{unsrt}

\runninghead{Visualising emotions} % A shortened article title to appear in the running head, leave this command empty for no running head

\footertext{\textit{PREPRINT} (2023)} % Text to appear in the footer, leave this command empty for no footer text

\setcounter{page}{1} % The page number of the first page, set this to a higher number if the article is to be part of an issue or larger work

%----------------------------------------------------------------------------------------
%	TITLE SECTION
%----------------------------------------------------------------------------------------

\title{Visualising emotions with linearly parameterized facial expressions} % Article title, use manual lines breaks (\\) to beautify the layout

% Authors are listed in a comma-separated list with superscript numbers indicating affiliations
% \thanks{} is used for any text that should be placed in a footnote on the first page, such as the corresponding author's email, journal acceptance dates, a copyright/license notice, keywords, etc
\author{%
	Paramoji
}

% Affiliations are output in the \date{} command
%\date{\footnotesize\textsuperscript{\textbf{1}}School of Chemistry, The University of Michigan\\ \textsuperscript{\textbf{2}}Physics Department, The University of Wisconsin\\ \textsuperscript{\textbf{3}}Biological Sciences Department, The University of Minnesota}

% COMMANDS
\usepackage{amsmath}
\usepackage{tabularx}



% Full-width abstract
\renewcommand{\maketitlehookd}{%
	\begin{abstract}
    \noindent
		The key challenge of affective computing is to translate a subjective emotional experience into an objective state.
Most recent advances in that field have relied on facial expressions as indicators for inner emotional states.

Two main approaches have emerged to accomplish this task: the dimensional and
the categorical approach. The former describes an emotion along a set of independent dimensions, such as valence, arousal,
and dominance. This is mathematically elegant, but suffers from lacking consensus on how
to represent more complex emotions. The categorical approach uses discrete emotional states
that are often defined by their corresponding facial expression. This approach has been
popularized through the widespread use of emojis that show such expressions and which humans can instantly recognize.
However, their categorical nature make it difficult to compute numerical calculations and statistical measures.

This paper seeks to bring the power of facial visualization to the dimensional approach.
It introduces a new dimensional model and a corresponding graphical representation.
According to five emotional dimensions (valence, arousal, dominance, contempt and control)
a unique visualization with strong emotional content is derived. Their comical style is free of references to gender, age or race.
Each such face could be an emoji, but they are parametric and allow all intermediate states to be represented.
To test their effectiveness I hired untrained crowd workers to rate the visualizations
along the five emotional dimensions. The results confirm that the encoded dimensions can indeed be
recorgnized intuitively and without further training.
	\end{abstract}
}

%----------------------------------------------------------------------------------------

\begin{document}

\maketitle % Output the title section

%----------------------------------------------------------------------------------------
%	ARTICLE CONTENTS
%----------------------------------------------------------------------------------------
\section{Introduction}

Humans have largely embraced the use of emoticons or emojis with facial expressions
to show emotions when communicating over a digital
channel. There is no doubt that these small faces fill a gap in informal written
communication that is hard to match with written words.
A vast amount of human generated content filled with such emotional cues
could have provided a unique opportunity for humans and computers to exchange emotional
information, as envisioned in the field of Affective Computing\cite{Picard97}.
However, this is not happening. There is no known application where humans and computers interact
with the meaningful use of emoji or other facial expressions.
Quite to the contrary, the lacking ability of showing appropriate emotional indicators
hampers the field of robotics, as the missing empathy makes robots seem less
intelligent and less trustworthy\cite{EmotionTowardsRobots, CHUAH2021102551}.
Visualizing an appropriate emotion with an appropriate facial expression seems to be a very difficult task for current algorithms.
What makes this emotional content so intuitive for humans, but yet so hard to process for computers?

The natural starting point for the analysis of digitized facial expressions are of course the standardized emojis.
Many attempts have been made to measure (or embed) these little icons with various
axes of emotional intensity\cite{SentimentOfEmoji, WhatDoesThisEmojiMean}.
Unfortunately, there is a lack of consensus among humans on how to evaluate different facial icons on a unified scale.
Cross cultural variations significantly impact how emojis are ranked on a single sentiment axis\cite{emoji_intensity}.
Things get more complicated when system borders are crossed.
The Unicode Consortium defines only rough graphical gidelines for many of the emotionally expressive symbols.
What is defined as a ``grinning face'' was found to have
at least 17 different graphical versions on various software platforms leading to
a severe misunderstanding on even basic attributes of sentiment ranging from positive to
negative\cite{emoji_interpretation}.
In an anatomy based approach the facial features in the emojis were studied\cite{emoji_facs}.
Inferring muscle activity and causing emotions proved difficult in the light graphical abstractions and tiny rendering differences.
The bottom line is that standardized emojis are
not suitable to encode nuanced and precisely defined emotions.

Outside the world of standardized emojis there are other sets of emoticons, i.e.
graphical facial expressions, to fill the need for fine grained scales of emotions.
Research has found them useful for the communication of qualias, such as pain and taste.
Two effective scales of faces were designed to capture levels of perceived pain within children and adults\cite{bakerwong, dvprs}.
It was found that visual scales of facial expressions produced more stable and consistent results in self-reported
pain and greatly improved treatment.
Likert scales are often used in surveys to assess agreement. Various designs of smiley faces were tested and
shown to impact engagement in children\cite{likertscale}.
A wider range of emotional experiences can be expressed if additional dimensions are allowed.
An emoji grid spans two dimensions, valence and arousal, with selected emojis along the axes\cite{emojigrid}.
It was found useful for reporting food elicited emotions.
Despite the additional depth of two dimensions the grid still lacks space for many emotions.
In short, these custom sets are consistenty defined and outperform words on their specific domain.
However, their expressive range is limited and does not nearly cover the variety of Unicode faces.

For a greater range of facial expressions there is active research into facial visualization techniques.
This includes 2D graphics\cite{parametric_face_facs, FACSGen} as well as physical 3D models\cite{FARAJ2021e00117}.
These techniques seek for realism in the displayed face and are loosing on the expressiveness of the exaggerated
emojis. Additionally, they mimic certain facial expressions without properly explaining their meaning in the digital realm.
Another --almost ancient-- technique for facial visualization uses line graphics
driven by high dimensional data\cite{chernoff}. Chernoff faces provide unique facial features controlled by 18 parameters covering the
sizes of eyes, mouth and ears as well as angles and distances between them.
Unfortunately, this does not include any of the subtle activations that are necessary for the display of deep emotions.
Chernoff faces come closest at being a predecessor, or next best solution, to the one I present in this paper.
Both share two key features: they use line art and an exaggerated comic style to visualize data and neither
has yet produced a success story in any practical application.
In the following sections I will do my best to let latter change.

\section{The parametric emoji}
The parametric emoji  (paramoji)
is a graphical icon that displays a human emotion from a set of continuous parameters.
The first question to answer is how an emotion can be mapped to numbers.
The idea of expressing a human emotion along measurable dimensions dates back to Wundt in 1896, who derived
his ideas mainly through self-inspection\cite{wundt}.
In a subsequent review \cite{tridimensional} the three dimensions of feeling were
translated as pleasantness, excitement and tension. Since then lots of different multi-dimensional models
for emotions were derived\cite{Cochrane2009-EDF, study5, ZHOU202038}.
This publication uses the dimensions ``valence'', ``arousal'', ``dominance'', ``contempt'' and ``control'',
partly inspired by previous emotional theories, party created to fit the specific expressability of a face.

As in all models it will be easy to understand the definition of each dimensions in isolation.
In combination it quickly becomes unclear how to interpret the resulting textual description.
To solve this problem the visualization in the form of a parametric emoji, a comic style
face icon, provides a unique and valid facial expression for each choice of parameters.
The face is made to be generic, without any reference to physiological properties, like age, gender or race.
Strength of the emotional expression and the algorithmic simplicity had top design priorities.
All graphical parameters depended linearly on the emotional dimensions, with the exception of contempt where
few bilinear factors drive the asymmetric look.

It would, of course, have been preferable to construct a scale that is linear in the expressed emotional intensity.
However, we would hardly be able to agree on a method for its measurement. Linearity in graphics is not just the second best option.
It is objective, as it is easy to check that the generating code essentially consists of a matrix vector multiplication.
It is at least something linear to which other systems can be calibrated. Finally, linearity creates the benefit that any detected smooth
transition from one state to another would also show smoothly and without artefacts when animated. A midpoint
in parameter space is guaranteed to also be the midpoint in its graphical representation.

\subsection{Valence and arousal}

The first two dimensions, known as ``valence'' and ``arousal'', are taken from the Circumplex model\cite{circumplex}.
The original publication describes a quantitative framework
to define all emotional states with two numerical parameters.
Valence describes the emotion along the axis of pleasure vs misery and approval vs disapproval. Arousal describes the axis
from aroused to sleepy, active to inactive or agitated to calm.
The claim that all emotions fall into a plane of only two dimensions might appear as quite a stretch.
However, this simplicity makes the model quite attractive and it became the basis of almost every emotional model.
Recent neurological evidence provides further assurance that these dimensions actually exist and correlate to
measurable brain activities. \cite{circumplex2}.

Figure \ref{fig:circumplex} shows the chosen visualization for the two dimensions.
The graphics is defined through the control points of Bezier lines with some rotation of the eyes' outlines.
All coordinates and angles are mapped linearly from the two input dimensions.
In this (still incomplete) model arousal controls eye and mouth sizes and valence controls curvature.

\begin{figure}[htbp!]
    \begin{center}
        \begin{tabular}{cc|ccc}
            && \multicolumn{3}{c}{valence $(v)$} \\[0.2em]
            && 0.0 & 0.5 & 1.0 \\ \hline
            &&&&\\[-0.8em]
            & \rotatebox{90}{\hspace{1.5em} 1.0} &
            \paramoji{0,100,100,50,0} &
            \paramoji{50,100,100,50,0} &
            \paramoji{100,100,100,50,0} \\
            \rotatebox{90}{\hspace{0.8em} arousal $(a)$}
            & \rotatebox{90}{\hspace{1.5em} 0.5} &
            \paramoji{0,50,50,50,0} &
            \paramoji{50,50,50,50,0} &
            \paramoji{100,50,50,50,0} \\
            & \rotatebox{90}{\hspace{1.5em} 0.0} &
            \paramoji{0,0,0,50,0} &
            \paramoji{50,0,0,50,0} &
            \paramoji{100,0,0,50,0}
        \end{tabular}
    \end{center}
    \caption{Facial expressions created with three different values for valence ($v$) and arousal ($a$).
    All displacements of graphical control points depend linearly on these two dimensions and can be freely interpolated. }
    \label{fig:circumplex}
\end{figure}

There is an obvious zero point for arousal when eyes and mouth stay totally shut.
On the other end of the spectrum there is no physical limit on how wide a comic eye can open.
Therefore, arousal values can only be compared against each other without an objective center.
For valence the natural center is a flat mouth and flat eyes.
Eyes and mouth bend upwards for pleasure.
This is consistent with a respective muscle movements and was already part of early emoticons
as \texttt{\^\_\^} and \texttt{:)}.
Eyes bend downwards for misery. This is derived from a drooping upper eyelid and the foreshortened view when the head is lowered.

Looking at facial expressions one can find the predicted emotions on the edges of the domain:
alarmed \paramojib{50,100,100,50,0},
sleepy \paramojib{50,0,0,50,0},
happy \paramojib{100,50,50,50,0} and
miserable \paramojib{0,50,50,50,0}.
The quadrants also comply with the Circumplex model, showing
contentment \paramojib{100,0,0,50,0},
excitement \paramojib{100,100,100,50,0},
distress \paramojib{0,100,100,50,0}
and depression\paramojib{0,0,0,50,0}.
The original publication postulates many more emotions that are supposed to sit tightly packed on a planar space.
However, fear vs anger vs tension vs disgust: these are fundamentally different emotions that can hardly be seen in close proximity to each other.
More dimensions are needed to separate them.

\subsection{Dominance}

This document uses the name ``dominance'' for the third dimension,
although it was first introduced as ``potency'' in  earlier publications.
The evaluation-potency-activity model was developed to find the main axes of meaning in nouns
by mapping them to adjectives through a data driven approach\cite{potency_osgood}.
Factor analysis revealed that two of the dimensions with most explanatory power were equivalent to valence and arousal.
The third of the three most important dimensions of meaning is described by the adjective scales of strong vs weak,
hard vs soft, brave vs cowardly, tenacious vs yielding.
Unfortunately, the authors did not find a good name for it.
The pleasure-arousal-dominance model fixed that problem \cite{dominance_mehrabian}. Its dimensions are
introduced as ``analogues'', because the results of this factor analysis basically replicated the dimensions with slightly changed names.
In this model anger, power and boldness dominate over anxiety, infatuation, fear and loneliness.

Irrespective of naming and definition details, the third dimension - dominance - shows anger and strength \paramojib{50,50,50,100,0}
on one end of the scale against fear and weakness \paramojib{50,50,50,0,0} on the other end.
The visual effects are lowered inner eyebrows, inwards rotated eyes, more prominent display of teeth and a sharper mouth contour.
For low dominance the visual effects are exactly mirrored, because all displacements are, again, linearly added.
Figure \ref{fig:dominance} shows three example paramojis with lowest (0.0) and highest (1.0) values of dominance.
Previous faces above are assumed to be at medium dominance (0.5). Obviously, the scale is continuous and all intermediate values could be interpolated.

\begin{figure}[htbp!]
    \begin{center}
        \begin{tabular}{l|ccc}
            \rotatebox{90}{\hspace{1ex}$d = 1.0$} &
            \paramoji{50,50,50,100,0} &
            \paramoji{20,20,20,100,0} &
            \paramoji{20,80,80,100,0}
            \\
            \rotatebox{90}{\hspace{1ex}$d = 0.0$} &
            \paramoji{50,50,50,0,0} &
            \paramoji{20,20,20,0,0} &
            \paramoji{20,80,80,0,0}
            \\[1ex]
            \hline\\[-1.8ex]
            $v$ & 0.5 & 0.2 & 0.2 \\
            $a$ & 0.5 & 0.2 & 0.8
        \end{tabular}
    \end{center}
    \caption{Three facial expressions plotted with high (1.0) and low (0.0) dominance.
    Although impossible to prove in print, all parameters can be combined arbitrarily.}
    \label{fig:dominance}
\end{figure}

The three parameters, valence, arousal and dominance, range from 0 to 1 and span the unit cube domain.
Four corners on the low valence side correspond to misery and disapproval. They already resemble some
possible versions of the basic emotions defined by Paul Ekman:
fear \paramojib{0,100,100,0,0},
anger \paramojib{0,100,100,100,0},
sadness \paramojib{0,0,0,0,0} and
disgust \paramojib{0,0,0,100,0}.
For positive valence we find pleasure and approval. Two universal emotions sit on the positive end of the cube:
happiness \paramojib{100,50,50,50,0} and
surprise \paramojib{100,100,100,50,0}.
The expressions in the corners of the positive side are probably not universal. I'll call them
simper \paramojib{100,0,0,0,0},
embarrassed \paramojib{100,100,100,0,0},
anticipating \paramojib{100,100,100,100,0}
and a forced smile \paramojib{100,0,0,100,0}.

\subsection{Contempt}

The fourth dimension shall be ``contempt'', derived from a discovery that this feeling of superiority is not a distinct emotion,
but combines differently with positive and negative feelings \cite{contempt_ekman}.
It is described by the adjectives scornful, haughty, smug, vain, or disdainful.
New findings also suggest that ``annoyance'' is an adequate label for this emotion\cite{contempt_new}.
Contempt is a feeling of moral superiority that varies in expression based on the amount of social support
one gets. Contempt against inferiors, the feeling of power can be a pleasurable experience.
Contempt against a superior can feel misarable. Because the expression depends on other dimensions it is a dimension
itself, rather than a specific spot in parameters space.

The only universally confirmed expression of contempt is the ``unilateral lip-tightening expression''\cite{contempt_ekman}.
A raised upper lip often shows in conjunction disgust. Tension in the cheek can also press the eyes into
squeezing. A squinting eye is not in itself a signal for contempt,
but a graphical choice to emphasizes its asymmetric nature and
often occurs in coincidence. All effects are limited to the left side of the face, an arbitrary choice,
because there is no known emotional meaning to chirality.

Figure \ref{fig:contempt} shows three emotions on two contempt levels. At least graphically,
contempt mixes perfectly with the previous dimensions. The left eye and the left side of the mouth have a bilinear
graphical dependence. Respective coordinates depend on $c$, but also on pairwise products with the other dimensions.
All combinations of parameters produce sometimes comical, but always valid facial expressions.

\begin{figure}[htbp!]
    \begin{center}
        \begin{tabular}{l|ccc}
            \rotatebox{90}{\hspace{1ex}$c = 1.0$} &
            \paramoji{86,19,19,68,100} &
            \paramoji{10,10,10,90,100} &
            \paramoji{38,72,72,41,100} \\
            \rotatebox{90}{\hspace{1ex}$c = 0.0$} &
            \paramoji{86,19,19,68,0} &
            \paramoji{10,10,10,90,0} &
            \paramoji{38,72,72,41,0} \\
            \hline\\[-1.8ex]
            $v$ & 0.9 & 0.1 & 0.4 \\
            $a$ & 0.2 & 0.1 & 0.7 \\
            $d$ & 0.7 & 0.9 & 0.4
        \end{tabular}
    \end{center}
    \caption{Random paramojis with and without contempt. The contempt
    dimension produces valid facial expressions when combined with any values in the previously defined dimensions.}
    \label{fig:contempt}
\end{figure}

\subsection{Control}

The final dimension of this model is ``control''. It describes the modulation of a
facial expression that is partly put under conscious control.
Hence, the interpretation of the last dimensions may partly be subject to cultural norms as
these demand certain interventions.
It ranges from controlled, modulated, suppressed \paramojib{50,100,0,50,0}
to uncontrolled, overwhelmed, expressed \paramojib{50,0,100,50,0}.
The visualization follows an hypothesis raised by Ekman and Friesen\cite{Ekman75}. They wrote: ``when a person is
controlling what is shown on his face, more effort is focused on managing what occurs in and around his mouth and
lips than in the area of the eyes/lids or brows/forehead.''  Although there might be limitations to that observation,
it is certainly good enough as a first approximation to the subject matter.

The major and easiest to control aspect of an emotion is its intensity. In visual terms, this is arousal ($a$). All
the graphical impact caused by $a$ can be separated into arousal shown in the eye region ($a_1$) and arousal
shown in the mouth region ($a_2$). The amount of excerted control ($k$) can be computed as the difference between
these two. For a controlled and modulated expression $k$ is large. For an uncontrolled, overwhelmed expression $k$ is low.
The dimensions $a_1$ and $a_2$ are mathemtically preferred, because they are defined on the Carthesian domain $[0,1]^5$.
The dimensions $a$ and $k$ have, at least theoretically, an emotional interpretation, but produce valid visualizations
only on a transformed domain.

Equation \eqref{eqn:control} defines how the dimension can be converted.
The complexity of the equations is commensurate with the psychological depth of this theoretical model. An empirical adjustment based on real human responses to the generated graphics is presented in a later section.

\begin{equation}
  \label{eqn:control}
  \begin{array}{rcl}
      a_1 &=& a + k/2 \\
      a_2 &=& a - k/2 \\
      k &=& a_1 - a_2 \\
%      a &=& (a_1 + a_2)/2
  \end{array}
\end{equation}

Figure \ref{fig:control} shows the visual effect of the control dimension on some random random faces.
It is interesting that even in the extremely overwhelmed versions one can still differentiate between
positive \paramojib{100,0,100,50,0}
and negative \paramojib{0,0,100,50,0} valence.
This is in contrast to real faces, where extreme pain and extreme joy cannot be distinguished from
the facial expression without context information or body language\cite{painjoy}.
Emoji style visualization are not just exaggerated versions of existing facial expressions.
They can also encode emotional information that go beyond what real faces show.

\begin{figure}[htbp!]
  \begin{center}
      \begin{tabular}{l|ccc}
          \rotatebox{90}{\hspace{1ex}$k = +1$} &
          \paramoji{90,100,0,50,0} &
          \paramoji{30,100,0,0,0} &
          \paramoji{10,100,0,80,70}
        \\
          \rotatebox{90}{\hspace{1ex}$k = -1$} &
          \paramoji{90,0,100,50,0} &
          \paramoji{30,0,100,0,0} &
          \paramoji{10,0,100,80,70}
          \\[1ex]
          \hline\\[-1.8ex]
          $v$ & 0.9 & 0.3 & 0.1 \\
          $a$ & 0.5 & 0.5 & 0.5 \\
          $d$ & 0.5 & 0.0 & 0.8 \\
          $c$ & 0.0 & 0.0 & 0.7
      \end{tabular}
  \end{center}
  \caption{Random emotions with minimal and maximal values of control ($k$) and $a=0.5$.
  This dimension reflects that humans tend to excert conscious control over the mouth more than the eyes.}
\label{fig:control}
\end{figure}

\section{Empirical evidence}

In the previous section I described my subjective choices for the paramoji's visual appearance.
In this section I will lay out my process to verify or adjust the chosen parameters based on empirical evidence.
In particular I wanted to know if untrained random participants would produce annotations of emotional dimensions
in consistency with their theoretical derivation.

\subsection{The experiment}

Figure \ref{fig:ui} shows the user interface that test participants had to navigate for the annotation of
emotional dimensions. Two faces were randomly drawn from a uniform distribution. The users' task was to order
them on the assigned emotional dimension. One had to be chosen as more intense and the the other consistently as
less intense. Alternatively, both faces could be labeled as equal. When the two choices matched
participants could continue to the next pair. No additional explanation was given after following a link
to this emotion classificion task.

\begin{figure}[ht!]
 \begin{center}
  \includegraphics[height=6cm]{gen/compare.pdf}
  \vspace*{-3ex}
 \end{center}
 \caption{
   User interface for untrained test participants. Random pairs of faces had to be ordered on
   one randomly assigned emotional dimension.
 }
 \label{fig:ui}
\end{figure}

Each participant was assigned to one emotional dimension, to avoid any cross influence caused by switching contexts.
Each dimension was described by adjectives for their negative and positive extreme. Synonyms were given to avoid
miscommunication. The following list shows the five possible tasks descriptions to which participants were assigned.

\begin{enumerate}
  \item Which face is more {\bf negative} (miserable, rejected); which is more {\bf positive} (pleasurable, accepted) ?
  \item Which face is more {\bf calm} (sleepy, relaxed, unimpressed); which is more {\bf aroused} (surprised, panicking, facing the unexpected) ?
  \item Which face is more {\bf submissive} (in doubt, passive, insecure); which is more {\bf dominant} (focused, strong, activated)?
  \item Which face is more {\bf uncontrolled} (overwhelmed, expressed); which is more {\bf controlled} (pretended, modulated, suppressed) ?
  \item Which face is more {\bf humble} (moderate, respectful); which is more {\bf contemptuous} (self-loving, vain, disdainful) ?
\end{enumerate}

\subsection{Participation}

Test participants were hired from MicroWorkers.com, a global market place for crowdsourced work.
A deep analysis of the workers' anatomy has been made previously and found to be still valid\cite{microworkers}.
The overall work force has a global bias towards India, Bangladesh and the Philippines with some participation
from the Western countries.

420 workers took part in this experiment and produced 13190 comparisons.
The pay was 1 Cent per entry, paid for batches of size 20 or 30. Some workers submitted two batches without extra reward.
The average time spent on one comparison was 10 seconds, the median 5 seconds.
It is likely that some workers have seen previous versions of the experiment, but never got any feedback on their performance.
The published data are the first results for a textual comparison task.

\subsection{Regression results}

The data were collected as five parameters for the Carthesian model and grouped by the
five dimensions that were rated. Additionally, a label was
stored to indicate whether it was seen as more, less or equally intense, yielding the rating numbers 1, 0 and 0.5.
Since all pairs were chosen fully at random
both sides of the comparison could be treated as independent samples. The average rating value of a
position in five-dimensional parameter space can be interpreted as the fraction of
paramojis that it supercedes, i.e. its quantile on the measured emotional scale.

A linear regression of the reported rating numbers onto the Carthesian dimensions reveals
how much the report is swayed by each
theoretical dimension. Table table \ref{tab:raw_data} shows the results.
The 1-norm column evaluates biggest possible change of the predicted quantile in two opposite corners
\paramojib{100,100,0,0,0} and
\paramojib{0,0,100,100,100}.
In absence of doubt or disagreement the least intense position should yield 0 and the most intense should yield 1.
As we can see the reported change in valence can be explained through the graphical transformations\footnote{
The value for valence is greater than the valid probability range from 0 to 1.
This is because the parameters were estimated based on random points within the domain.
The corner points were not among the queried samples and thus the results here are extrapolated values}.
It is also interesting that just over half the effect on valence originates from the iconic
mouth curvature controlled by $v$,
which many icon sets depend on solely. Pleasure is clearly more than a mouth shape.

The next thing to observe in table \ref{tab:raw_data} is that most rows are dominated by the $v$ column.
With exception of dominance
there is a high degree of colinearity between all rows. An aroused or contemptuous face was
likely also seen as negative. Control was
seen as positive. The describing adjectives were not understood as independent.
Since we are interested in orthogonal dimensions, we must map these results onto independent axes.

\begin{table}[ht!]
  \begin{center}
    \begin{tabular}{lrrrrrr}
      {} &     $v$ &     $a_1$ &     $a_2$ &     $d$ &     $c$ &  1-norm \\
      \midrule
      Valence $(v^e)$   &  0.54 & 0.07 & -0.06 & -0.21 & -0.16 & 1.03 \\
      Arousal           & -0.29 & 0.01 &  0.12 &  0.16 &  0.06 & 0.63 \\
      Dominance         & -0.05 & 0.01 & -0.01 &  0.19 &  0.01 & 0.29 \\
      Control           &  0.28 & 0.10 & -0.16 & -0.10 & -0.10 & 0.75 \\
      Contempt          & -0.24 & 0.02 &  0.03 &  0.20 &  0.08 & 0.58 \\
      \midrule
    \end{tabular}
  \end{center}
  \caption{Regression results of the likelihood of an emotional attribute being reported as more intense onto the Carthesian visualization parameters.}
  \label{tab:raw_data}
\end{table}

\subsection{Orthogonalization}

Apparently, respondents did not understand the meaning of the questions as orthogonal.
A better question to ask would have been along the line of “who looks more aroused,
without considering the effect of negativity?” However such questions would become quite
confusing when adding more dimensions.

The Gram-Schmidt method is a standard procedure to calculate such an orthogonalization mathematically.
Equation (\ref{eqn:gramschmidt}) shows the calculation. The vectors v are the linear dependent
input shown in the rows of table \ref{tab:raw_data}. The orthogonalized output u is seen in the rows
of table \ref{tab:ortho_data}.

\begin{equation}
  \mbox{u}_i
  \,=\,
  \mbox{v}_i - \sum\limits^{i-1}_{j=1}
    \frac{\langle
      \mbox{v}_i, \mbox{u}_j
    \rangle}
    {\langle
      \mbox{u}_j, \mbox{u}_j
    \rangle}
    \mbox{u}_j
  \label{eqn:gramschmidt}
\end{equation}

The Gram-Schmidt method depends on an order of the input and yields slightly different results when rearranged.
Some thought should be put into this choice to ensure correct interpretation of the results.
The first dimension should clearly be the valence direction. It is the one most decisively identified by participants. It is
also the most commonly agreed dimension in literature with some research using it as the only measured variable.
Table \ref{tab:ortho_data} shows the results starting with valence and then ordered by decreasing 1-norm of the remaining vector.

\begin{table}[ht!]
  \begin{center}
    \begin{tabular}{lrrrrrr}
      {} &         $v$ &     $a_1$ &     $a_2$ &     $d$ &     $c$ &       1-norm \\
      \midrule
      Valence $(v^e)$   &  0.54 & 0.07 & -0.06 & -0.21 & -0.16 & 1.03 \\
      Dominance $(d^e)$ &  0.05 & 0.03 & -0.02 &  0.15 & -0.02 & 0.27 \\
      Control $(k^e)$   & -0.03 & 0.05 & -0.12 & -0.02 & -0.01 & 0.24 \\
      Arousal $(a^e)$   & -0.02 & 0.06 &  0.04 & -0.00 & -0.03 & 0.15 \\
      Contempt $(c^e)$  &  0.01 & 0.01 &  0.00 & -0.00 &  0.03 & 0.05 \\
      \midrule
    \end{tabular}
  \end{center}
  \caption{Orthogonalized results. This table shows the independent impact of each dimension after
  applying Gram-Schmidt orthogonalization.}
  \label{tab:ortho_data}
\end{table}

\subsection{Orthogonalized results}

The orthogonalization started with the emperical measurement for valence $(v^e)$. As the first vector it was
not changed by the process.
Figure \ref{fig:emp_valence} shows the result as a graphical transformation, which
can be interpreted to sway the perceived sentiment from negative to positiv most efficiently, i.e. with
minimal graphical displacements.
For the visualization it is applied to a neutral face
\paramojib{50,50,50,50,25}, which is chosen at the parameters $(0.5, 0.5, 0.5, 0.5, 0.25)$.
This is a purely arbitrary choice. It just so happens to be in the center of the domain with a bias to
the non-contemptuous symmetric expression. Any other parameters could have been chosen as the center, because the
emotional effect of the graphical transformations were measured across the entire domain and not at a single point,
or along a single axis. From the chosen center all parameters are moved in proportion to vector $v^e$ until they reach the boundaries of the drawable space in the unit cube.

\begin{figure}[ht!]
  \begin{center}
    \begin{tabular}{cccc}
        \paramoji{0,43,56,70,40}
        \paramoji{33,48,55,57,30}
        \paramoji{67,52,50,43,20}
        \paramoji{100,57,52,36,10}
    \end{tabular}
  \end{center}
  \caption{Emperical data suggest that this is the most effective graphical change to sway reported impressions from negative to positive.}
  \label{fig:emp_valence}
\end{figure}

The second strongest orthogonal dimension is dominance $(d^e)$.
Figure \ref{fig:emp_dominance} shows the faces on the dominance dimension after removing all the colinear projections.
Dominance is almost naturally understood as an independent dimension. Little correction is necessary, but
the mouth on the dominant side shows a light smile, which would be positive. However, the frowning eyebrows are considered negative.
According to this analysis the effects cancel each other such that there is no linear measurable
valence change on this scale.

\begin{figure}[ht!]
  \begin{center}
    \begin{tabular}{cccc}
        \paramoji{34,41,57,0,30}
        \paramoji{45,47,52,33,27}
        \paramoji{55,53,48,67,23}
        \paramoji{66,59,43,100,20}
    \end{tabular}
  \end{center}
  \caption{Emperical evidence for dominance confirms the eye brow angle as the driving factor.}
  \label{fig:emp_dominance}
\end{figure}

The third strongest dimension is control $(k^e)$ in figure \ref{fig:emp_control}.
Although control is correlated to positive sentiment in the raw data, the orthogonalization process seems to slightly overcompensate.
Anyways, the main driver is confirmed as the difference between the opening of eyes and mouth.

\begin{figure}[ht!]
  \begin{center}
    \begin{tabular}{cccc}
        \paramoji{63,28,100,57,30}
        \paramoji{54,43,67,52,27}
        \paramoji{46,57,33,48,23}
        \paramoji{37,72,0,43,20}
    \end{tabular}
  \end{center}
  \caption{Emperical reports for control show the scale from throwing a tantrum to the proverbial biting of one's tongue.}
  \label{fig:emp_control}
\end{figure}

The fourth empirically confirmed dimension is arousal $(a^e)$ in figure \ref{fig:emp_arousal}.
After freeing the reported numbers from the effects of negativity and dominance, the expected
synchroneous opening of eyes and mouth emerges. This effect adds up to 0.15 to the probability of
being more aroused than a random face. This is a low value given its status as the number two in most models.
Apparently there is some limit to the extent to which a static image can convey the impression of arousal.
This dimension would be expected to coincide with sudden movements and rapid changes that cannot be shown in print.
Nevertheless, the effect is measurable and follows expectation.

\begin{figure}[ht!]
  \begin{center}
    \begin{tabular}{cccc}
        \paramoji{62,0,22,54,49}
        \paramoji{54,33,41,51,33}
        \paramoji{46,67,59,49,17}
        \paramoji{38,100,78,46,1}
    \end{tabular}
  \end{center}
  \caption{Empirical results confirm that synchroneously opening eyes and mouth make a face look aroused.}
  \label{fig:emp_arousal}
\end{figure}

The final dimension is basically whatever remains as independent from previous observations.
Figure \ref{fig:emp_contempt} visualizes this orthogonal remnant.
This confirms our expectations of a contemptuous look, although it only
explains 0.05 of the probability. Despite being barely measurable
we have no reason to doubt the correctness of this result. The asymmetry is a unique facial expression
and well documented, though possibly not widely understood.
In addition to the asymmetry in mouth and eyes there is an even further enlarged right eye.
It can thus be assumed that the asymmetry in
the eyes helps in making this emotion more recorgnizable.

\begin{figure}[ht!]
  \begin{center}
    \begin{tabular}{cccc}
        \paramoji{44,38,48,51,0}
        \paramoji{52,54,51,50,33}
        \paramoji{59,70,53,48,67}
        \paramoji{67,86,55,47,100}
    \end{tabular}
  \end{center}
  \caption{The contempt dimension is the remnant that cannot be explained by previous dimensions.
    There is no doubt; this asymmetric look is indeed contempt.}
  \label{fig:emp_contempt}
\end{figure}

\subsection{Transformation}

Formula (\ref{eqn:transform}) summarizes the transformation from empirical coordinates into Cartesian ones for plotting.
They are ordered by decreasing effect on the reported emotional impression.
Each entry changes the impact of its emotional component without affecting the previous ones.
The function maps the default point $(0.5, 0.5, 0.5, 0.5, 0.25)$ onto itself.

\begin{equation}
  \left(
    \begin{array}{c}
      v \\ a_1 \\ a_2 \\ d \\ c
    \end{array}
    \right) =
    \left(
      \begin{array}{ccccc}
        +1.00 & +0.32 & -0.25 & -0.24 & +0.22 \\
        +0.13 & +0.18 & +0.44 & +1.00 & +0.48 \\
        -0.11 & -0.13 & -1.00 & +0.56 & +0.07 \\
        -0.39 & +1.00 & -0.14 & -0.07 & -0.04 \\
        -0.30 & -0.10 & -0.09 & -0.47 & +1.00 \\
      \end{array}
    \right) \left(
    \begin{array}{c}
      v^e \\ d^e \\ k^e \\ a^e \\ c^e
    \end{array}
    \right) + \left(
    \begin{array}{c}
      +0.03 \\ -0.50 \\ +0.83 \\ +0.31 \\ +0.48
    \end{array}
    \right)
    \label{eqn:transform}
\end{equation}

The entries in the matrix are scaled to map each axis through the default point into a valid line within the unit cube.
This ensures that each individual axis can always be visualized with the appended plotting function.
Off-axis points in the empirical domain may be mapped outside the unit cube of the Cartesian model.
This is not such a big problem. As long as $a_1$ and $a_2$ are positive and keep the eyes at least as open as shut,
the face can be extrapolated and
recorgnized with some acceptance for artistic abstractions. Figure \ref{fig:extra} shows some results.


\section{Conclusion}
This article introduced a parametric facial expression that can be parametrized to
display a wide variety of emotions. Unlike other graphical methods, each parameter
has a distinct emotional interpretation. Empirical evidence has been provided that
the emotional teint of each axis can be extracted by untrained test participants. Their rating
of the resulting faces along the scales of valence, arousal, dominance, contempt und
control reproduced the input parameters with minor adjustments.

It has become apparent that little facial expressions can add an emotional note
that is hard to match in written text.
Emojis have filled that role in human to human communication. They owe their
initial success largely to facial expressions with emotional cues.
The presented visualizations are similarly suited for the exchange of emotional information.
They synthesize a comparable range of facial expressions, hence called, parametric emoji or paramoji for short.
However paramojis are much more suitable for computer processing,
as questions for similarity, intensity and sentiment are straight forward to resolve.
Stasticial measures like averages are supported in a meaningful manner.

It is known, that cultural norms demand varying degrees of self control with
regards to the display of emotions in public. Consequently, facial expressions
are interpreted according to these expections. When mapping emotions to adjectives
language skills and educational background also play a role. A method was presented
to adjust the emotion parameters to a specific target audience, crowd workers
in this case. Only linear adjustments were considered, but there is reason to believe
that non linear effects are small. The graphical transformations in each
of the five dimensions is extremely smooth, mostly linear, in all places.
When animated, all transitions always looks smooth and straight, making them
ideal for visualizations in real time or for replay.

\section*{Code}

The source code of the paramoji was difficult to construct, but is easy to read.
A sparse matrix contains all values that define the transformation.
There is a row for each free parameter
in the drawing instruction and a column
for each entry in the transformed input vector $(1, 2v-1, a_1, a_2, 2d-1, c, ca_1, cv, ca_2)$.
The dot product is computed between each matrix row and the input vector,
which yields the values for all the question marks in the template string.
The output format is SVG, a standardized description language for
vector graphics. Both, the JavaScript code and the SVG instructions, can be interpreted by any contemporary
web browser.\footnote{A demo of this code is currently deployed at \url{https://paramoji.org}}

{\small
\begin{verbatim}
function paramoji(v, a1, a2, d, c) {
 const matrix=[[,,,,,.5],[70,-2,,-8,-1],[-9,-1,-2],[37,-5,,,2,4],[68,-1,-2,-7,-2,
  -6,-4,-1,3],[29,-2,,-3,2,-9],[76,-1,,-8,,-17,-1,,3],[25,-4,,-6,2,-7],[77,-10,,-
  2],[65,4,,,-2,-.5],[68,-1,-2,-7,-2],[71,2,,3,-2],[76,-1,,-8],[75,4,,6,-2],[77,-
  10,,-2],[80,5,,6,-6,,,,-2],[66,2,,-5,-1,,,,-2],[14,1,,9,-1],[,-12,,-2,,-3,,6,-2
  ],[10,1,,11,-2],[,-2,,11,,,,,-9],[6,,,6,-2],[,1,,12,,,,,-5],[,1,,12,,,,,-4],[-
  10,-1,,-11,2],[,-2,,11,,,,,-2],[-14,-1,,-9,1],[,-12,,-2,,-3,,6,2],[-10,-1,,-11,
  2],[,-2,,-11,,,,,-5],[-6,,,-6,2],[,1,,-12,,,,,-3],[,1,,-12,,,,,-1],[10,1,,11,-2
  ],[,-2,,-11,,,,,-2],[14,1,,9,-1],[,-12,,-2,,-3,,6,-2],[,,,,,.4,,.15,-.25],[50,,
  ,,,-12,,,6],[81,2,,,-2,-2,,,1],[51,-2,-10],[,22,,,-26],[,,-11],[-8,,-8],[,,-4],
  [-8,,-8],[,,2],[-8,,-3],[,,9],[-4,,1],[,-6,13],[5,,12],[,-9,11],[11,,8],[5,,12]
  ,[,-9,-11],[-4,,5],[,-6,-13],[,,-11],[,,,,-1],[,-2.5,,,-1.5],[.5,,5],[5,1,-4,,8
  ],[2,-9,-17,,5],[,-22,,,26,-10,,20],[,,-11,,,,10],[8,,8,,,,-6],[,,-4,,,,4],[8,,
  8,,,,-6],[,,2,,,,-2],[8,,3,,,,-2],[,,9,,,,-9],[4,,-1],[,-6,13,,,-3,-12,6],[-5,,
  -12],[,-9,11,,,-5,-8,10],[-11,,-8],[-5,,-12],[,-9,-11,,,-5,9,10],[4,,-5],[,-6,-
  13,,,-3,12,6],[,,-11,,,,10],[,,,,2],[2,1,-10,,-6,3],[2,-9,-17,,4,20],[59,2,1,,-
  6],[40,-8,-19,,12],[71,2,3,,-5],[40,-11,-21,,3],[79,1,6,,-4],[40,-3,-18,,-3],[2
  ,-.5,1,,.5],[62,,-4,,-2],[33,-5,-16,,7],[56,2,,,-3],[39,-7,-22,,9],[56,1,,,-1],
  [41,-8,-20,,10],[,,,,,1,3,2]]
 const template='<svg width="100%" height="100%" viewBox="0 0 100 100"><defs><c'+
  'lipPath id="c-eye-r"><use href="#eye-r"/></clipPath><clipPath id="c-eye-l"><'+
  'use href="#eye-l"/></clipPath><clipPath id="c-lips"><use href="#lips"/></cli'+
  'pPath></defs><path transform="matrix(1,?,0,1,53,?)" d="M-4,0q-2,-2 -4,0M-1,0'+
  'Q3,-2 4,-1T6,-3 4,?" fill="none" stroke="black"/><path d="M?,?Q?,? ?,?" fill'+
  '="none" stroke="black"/><path d="M?,?Q?,? ?,?" fill="none" stroke="black"/><'+
  'g clip-path="url(#c-lips)"><rect height="100" width="100"/><ellipse cx="50" '+
  'cy="91" rx="15" ry="10" fill="#800f08"/><g transform="translate(0,?)"><use h'+
  'ref="#tooth" x="-14"/><use href="#tooth" x="-7"/><use href="#tooth"/><use hr'+
  'ef="#tooth" x="7"/></g><g transform="translate(0,?)"><use href="#tooth" tran'+
  'sform="matrix(1,.14,0,1,-14,-8)"/><use href="#tooth" x="-7"/><rect id="tooth'+
  '" x="50.5" rx="2" ry="1" height="15" fill="white" width="6" stroke="black" s'+
  'troke-width=".5"/><use href="#tooth" transform="matrix(1,-.14,0,1,7,7)"/></g'+
  '></g><path id="lips" d="M?,?C?,? ?,? 0,?S?,? ?,?C?,? ?,? 0,?S?,? ?,?Z" trans'+
  'form="matrix(1,?,0,1,?,?)" fill="none" stroke="black"/><g transform="transla'+
  'te(68,?)"><path id="eye-r" transform="rotate(?)" d="M-8,?C?,? ?,? ?,? ?,? ?,'+
  '? ?,0 ?,? ?,? -8,?Z" fill="white" stroke="black"/><g clip-path="url(#c-eye-r'+
  ')"><g id="lens" transform="translate(?,?)"><circle r="5" fill="#9d4922"/><ci'+
  'rcle r="?"/><circle r="1" cx="-2.5" cy="-1.5" fill="white"/></g><path d="M-2'+
  '5,-40H20V0Q?,? -15,-2Z" opacity=".25"/></g><g transform="translate(-36,0)"><'+
  'path id="eye-l" transform="rotate(?)" d="M8,?C?,? ?,? ?,? ?,? ?,? ?,0 ?,? ?,'+
  '? 8,?Z" fill="white" stroke="black"/><g clip-path="url(#c-eye-l)"><use href='+
  '"#lens" x="?"/><path d="M-50,-40H15V-2Q?,? -20,0Z" opacity=".25"/></g></g></'+
  'g><g id="brow-r"><path d="M?,?Q?,? ?,?" fill="none" stroke="black" stroke-wi'+
  'dth="?"/><path d="M?,?Q?,? ?,?" fill="none" stroke="black"/></g><use href="#'+
  'brow-r" transform="matrix(-1,0,0,1,100,?)"/></svg>'
 const dotprod = (X,Y) => X.reduce((a, b, i) => a + b*(Y[i] || 0), 0)
 const vec = [1, 2*v-1, a1, a2, 2*d-1, c, c*a1, c*v, c*a2]
 return template.replace(/\?/g, () => dotprod(vec, matrix.shift()))
}
\end{verbatim}
}
\vspace{1mm}

\begin{figure}[ht!]
  \begin{center}
    \begin{tabular}{cccc}
        \paramoji{105,138,95,106,-60}
        \paramoji{166,53,130,105,2}
        \paramoji{-55,124,126,131,-9}
        \paramoji{-58,180,0,-25,-9}
    \end{tabular}
  \end{center}
  \caption{Extrapolations of the facial features outside the valid range: recognizable, but increasingly grotesque.
  \texttt{\^{}\^}
  }
  \label{fig:extra}
\end{figure}


%----------------------------------------------------------------------------------------
%	 REFERENCES
%----------------------------------------------------------------------------------------

%\printbibliography % Output the bibliography
\bibliography{paramoji}

%----------------------------------------------------------------------------------------

\end{document}
